\subsection{辅助生殖各阶段的性生活指引：从降调到验孕}

进入辅助生殖（IUI / IVF）周期的夫妇，最常问却最少得到系统回答的问题之一是："这期间还能同房吗？"答案是：\textbf{绝大多数阶段可以，个别阶段需要暂停}——需要暂停的时间加起来通常不超过几周。本节按治疗时序给出一般规律；任何阶段都以生殖中心的个体化医嘱为准。

\vspace{0.5em}
\noindent\textbf{降调与促排卵阶段}：

\begin{itemize}
	\item \textbf{降调期间可以同房，但必须严格避孕}：降调与促排药物若在意外妊娠中继续暴露，可能干扰早期胚胎发育——本阶段同房请全程使用安全套，不要用"安全期法"去赌
	\item \textbf{促排中后期防卵巢扭转}：多卵泡发育使卵巢显著增大，悬韧带承受额外张力——剧烈的腹部撞击、深插入体位下猛然变换姿势，可能诱发\textbf{卵巢扭转}（突发一侧下腹剧痛、恶心呕吐，需急诊处理，参见前文风险部分）。此阶段建议温和体位、避免腹部受压
	\item \textbf{腹胀是身体在说话}：OHSS 高危者（年轻、瘦小、多囊卵巢综合征、AMH 偏高、预期取卵数多）腹胀加重时，欲望下降再正常不过——身体此时需要的不是"坚持完成任务"
\end{itemize}

\vspace{0.5em}
\noindent\textbf{人工授精（IUI）周期}：

\begin{itemize}
	\item \textbf{男方取精前禁欲 2--7 天}：过短或过长都会影响精液参数，按化验单要求执行
	\item \textbf{IUI 操作后}：当日避免性生活；无出血腹痛者多可自次日恢复温和性生活，部分中心建议 24--48 小时后——遵本中心医嘱；自然周期"指导同房"者按医嘱安排时机
\end{itemize}

\vspace{0.5em}
\noindent\textbf{取卵前后}：

\begin{itemize}
	\item \textbf{取卵前 1--2 天起暂停同房}：经阴道穿刺操作需要相对清洁的阴道环境，近期性生活增加操作区感染与出血的风险
	\item \textbf{取卵后 1--2 周内避免同房、盆浴与游泳}：穿刺针道需要时间闭合；期间出现阴道出血、腹痛加剧或发热，及时联系中心
\end{itemize}

\vspace{0.5em}
\noindent\textbf{胚胎移植后到验孕}：

\begin{itemize}
	\item \textbf{没有证据表明移植后短期同房会"冲掉胚胎"}：自然受孕中，着床期的性生活并不增加流产率；但多数中心仍建议移植后至验孕前暂停——既减少出血诱因，也避免"万一没着床是不是怪那一次"的终身自责
	\item \textbf{亲密不必等于插入}：两周等待期里，拥抱、依偎、按摩、口与手的亲昵（无不适前提下）都是联结——很多夫妇事后回忆，这段"只许温柔"的时期反而成了难得的慢时光
\end{itemize}

\vspace{0.5em}
\noindent\textbf{验孕之后}：

\begin{itemize}
	\item \textbf{阳性}：OHSS 未消退、阴道出血或腹痛期间禁止性生活；确认宫内活胎、症状消退后可逐步恢复（妊娠期性生活的一般原则见孕期相关章节）
	\item \textbf{阴性}：月经来潮、身体恢复后即可恢复正常性生活。不少夫妇在失败周期后需要一个"把生育与性分开"的缓冲期——这是心理修复需要，不是功能问题；若持续回避数月以上、开始影响关系，建议寻求心理支持（见本章心理与伦理部分）
\end{itemize}

\begin{tcolorbox}[colback=pink!5!white,colframe=red!50!black,title=贴心小叮咛]
取卵日男方"现场取精"紧张到取不出来的情况并不少见——这不是能力问题，是表现焦虑。预防很简单：提前与中心沟通\textbf{冻存一份精液备份}；有备份的人当天反而最放松，多数用不上它。
\end{tcolorbox}

本节是一般规律。每个方案的用药、卵巢反应与内膜准备方式不同——\textbf{任何阶段拿不准时，直接问主管医生}：这是生殖中心每天都会被问到的正常问题，不是"不好意思开口"的那类。
